% Capítulo 3 — Repaso de Álgebra del CBC con la mirada de Clavijo
\chapter{Repaso de Álgebra del CBC con la mirada de Clavijo}\label{cap:03}

\section{La idea}

% --- 03-cbc-nivel1 (redactado 2026-10-18) ---
\subsection{Lo que trajimos del CBC y qué le cambia la mirada}\label{sec:cbc-mirada}

En Álgebra del CBC cada objeto llegaba con su regla de cálculo: el producto escalar se
hacía multiplicando componente a componente y sumando, el determinante se desarrollaba por
una fila, el sistema se triangulaba con Gauss, la inversa salía de pivotear contra la
identidad. Todas esas reglas siguen valiendo y no las vamos a reemplazar. Lo que cambia es
la pregunta que hacemos antes de calcular: qué se mueve, qué queda quieto, qué se conserva,
con qué se mide. Cada cuenta del CBC tiene un significado en el terreno, y la materia
consiste en leerlo. El profesor lo planteó desde la primera clase: ``Transformaciones y
Operaciones (por ahora definidas coloquialmente) son dos conceptos que van atravesando la
materia'' (resumen 2025, clase 1). No agregamos álgebra nueva; cambiamos la pregunta.

La tabla recorre los temas del CBC en el orden en que suelen aparecer y dice, en una línea,
cómo los miramos acá y en qué capítulo vuelven.

{\small
\begin{tabularx}{\textwidth}{>{\raggedright\arraybackslash}p{30mm} Y >{\raggedright\arraybackslash}p{17mm}}
\toprule
\textbf{Lo que sabíamos (CBC)} & \textbf{Cómo lo miramos acá} & \textbf{Capítulo}\\
\midrule
Vector y punto: una flecha o una terna de números &
El punto no es su columna de números: el terreno es un espacio afín, hay mojones y
desplazamientos entre mojones; la columna aparece recién cuando elegimos origen y ejes. &
cap.~\ref{cap:04}, cap.~\ref{cap:06}\\
Producto escalar \(u\cdot v\) &
Dos cosas pegadas: fila por columna es una forma lineal, una regla de medir que lee un
punto sin necesidad de ángulos (\(\ell^Tp=0\) es incidencia: el punto está sobre la
recta); largos y ángulos son otra cosa, la métrica. &
cap.~\ref{cap:05}, cap.~\ref{cap:08}\\
Producto vectorial; determinante \(2\times2\) y \(3\times3\) &
Área o volumen con signo; el producto cuña es el ``detector de colinealidad'' (resumen
2025, clase 7): vale cero exactamente cuando los vectores están alineados. &
cap.~\ref{cap:08}\\
Recta y plano: \(ax+by=c\), \(ax+by+cz=d\) &
La implícita es una forma lineal \((a\ \ b)\) leída en el nivel \(c\): una máquina de medir
cuyas curvas de nivel son las rectas paralelas del mazo. La paramétrica \(P+t\,v\) es otra
cosa: la trayectoria de una traslación. &
cap.~\ref{cap:07}, cap.~\ref{cap:08}\\
Sistema lineal y Gauss &
Cada ecuación es una recta o un plano (el producto de matriz por vector ``parecía estar
formado por ecuaciones de rectas'', resumen 2025, clase 4); resolver \(Ax=b\) es encontrar
la preimagen de \(b\) bajo la transformación \(A\). &
cap.~\ref{cap:08}, cap.~\ref{cap:09}\\
Matriz y producto de matrices &
Una matriz es una transformación, no una tabla: ``las matrices y las transformaciones
lineales tienen una correspondencia 1 a 1'', y ``las Transformaciones se identifican con
las trayectorias que generan (por aplicación repetida)'' (resumen 2025, clase 1). El
producto es la composición. &
cap.~\ref{cap:04}, cap.~\ref{cap:08}\\
Inversa \(A^{-1}\) &
Deshacer el movimiento. La inversa se dice respecto de una operación: le llamamos a algo
``inversa con respecto a \ldots'' (resumen 2025, clase 10). Es el opuesto para la suma, el
recíproco para el producto, la transformación que devuelve cada punto a su lugar para la
composición. &
cap.~\ref{cap:08}, cap.~\ref{cap:09}\\
Cambio de base &
La misma transformación vista con otra regla: cambiar la base de su matriz es ``componer
esta transformación con las transformaciones de cambio de base que correspondan'' (resumen
2025, clase 8), \(B^{-1}AB\). &
cap.~\ref{cap:08}\\
Transformación lineal &
Un movimiento del espacio que deja quieto el origen y respeta sumas y escalas; el molde
para afinidades, proyectivas y, en lo chico, para cualquier transformación derivable. &
cap.~\ref{cap:06}, cap.~\ref{cap:08}\\
Autovalores y autovectores &
Dirección que cae sobre sí misma; en coordenadas homogéneas, punto fijo. Si la matriz tiene
tantos ``subespacios unidimensionales invariantes'' como la dimensión (resumen 2025,
clase 7), mirada desde ellos es diagonal. &
cap.~\ref{cap:07}, cap.~\ref{cap:08}\\
Norma y distancia &
La pone el producto interno, no el espacio: ``no existe una base vectorial privilegiada de
\(\R^n\)'' (resumen 2025, clase 6) salvo la que la métrica elige; sobre una superficie la
métrica cambia de punto a punto. &
cap.~\ref{cap:05}, cap.~\ref{cap:11}\\
\bottomrule
\end{tabularx}
}

\subsubsection*{Los cuatro cambios de chip, en \(\R\) y en \(\R^2\)}

Las cuatro maneras de soltar el reflejo del CBC que usa este curso (las unidades
interactivas las llaman ``cambios de chip'') las desarrollamos, como ladrillos, en la
sección~\ref{sec:desaprender}. Acá alcanza con un ejemplo mínimo de cada una, primero sobre
la recta y después en el plano.

\paragraph{1. No hay flechas flotando: todo vector es un punto o un desplazamiento entre
mojones.}
Sobre una traza, dos mojones están en las progresivas \(P=3\) y \(Q=7\), medidas desde el
mojón \(O\). Si medimos desde otro mojón \(O'\), ubicado en la progresiva \(1\), leemos
\(P'=2\) y \(Q'=6\). La diferencia no cambia, \(Q-P=Q'-P'=4\): es un desplazamiento. La suma
sí: \(P+Q=10\) señala la progresiva \(10\), y \(P'+Q'=8\) señala la progresiva \(8+1=9\).
Sumar puntos depende del origen; restarlos, o promediarlos, no
(\(\tfrac{P+Q}{2}=5\) y \(\tfrac{P'+Q'}{2}+1=5\)). En el plano pasa lo mismo: con
\(P=(1,2)\), \(Q=(4,6)\) y el origen corrido a \(O'=(1,1)\),
\[
  Q-P=(3,4)=Q'-P',\qquad P+Q=(5,8)\neq(3,6)+(1,1)=(4,7).
\]
El hecho geométrico: la flecha \((3,4)\) mide \(5\) desde cualquier origen; la columna
\((5,8)\) no es ningún lugar del terreno.

\paragraph{2. El gradiente no es una flecha, es una regla.}
Una forma lineal en \(\R\) es \(\omega(x)=3x\). Si pasamos de metros a medios metros, la
coordenada del punto se duplica, \(x'=2x\), y para que la lectura no cambie el coeficiente
de la regla se divide por dos: \(\omega'(x')=\tfrac32\,x'=3x\). Los puntos y las reglas se
transforman al revés. En el plano, la recta \(2x+y=4\) es la regla \(\omega=(2\ \ 1)\)
leída en el nivel \(4\); el punto \((1,2)\) está sobre ella porque \(2\cdot1+2=4\).
Estiramos el eje \(x\) con \(A=\mathrm{diag}(2,1)\): el punto va a \((2,2)\) y la regla va a
\(\omega A^{-1}=(1\ \ 1)\), que lo sigue leyendo en \(4\). Si en cambio movemos \((2,1)\)
como si fuera una flecha, obtenemos \((4,1)\), y la ``recta'' \(4x+y\) lee \(10\) en el
punto movido: ya no pasa por él. El hecho geométrico: una regla se dibuja como un mazo de
rectas de nivel equiespaciadas, y si el plano se estira, las marcas se juntan.

\paragraph{3. Las coordenadas son etiquetas, no metros.}
Si en la recta cada unidad de la grilla son \(2\) metros, la métrica es \(g=4\) y un salto
de \(3\) unidades mide \(\sqrt{g}\cdot3=6\) metros. En un plano dibujado con el Este a
\(2\) metros por unidad y el Norte a \(1\), la métrica es \(G=\mathrm{diag}(4,1)\). Los
vectores \(u=(1,1)\) y \(v=(1,-1)\) son perpendiculares en el papel, \(u^Tv=0\), pero en el
terreno
\[
  u^TGv=4-1=3,\qquad \|u\|_G=\|v\|_G=\sqrt5,\qquad \cos\alpha=\tfrac35,\quad
  \alpha\approx53{,}13^\circ .
\]
El hecho geométrico: ``perpendicular'' y ``largo'' son preguntas a la métrica, no a la
grilla.

\paragraph{4. Álgebra y análisis son la misma materia.}
En \(\R\), cerca de \(x=3\), \(f(x)=x^2\) se comporta como la multiplicación por \(6\):
\(f(3{,}01)=9{,}0601\approx9+6\cdot0{,}01\). La derivada es una matriz de \(1\times1\) que
escala los pasitos. En el plano, las coordenadas polares \((r,\theta)\mapsto(r\cos\theta,
r\sin\theta)\) tienen en \((2,0)\) la jacobiana \(J=\mathrm{diag}(1,2)\): un pasito
\(d\theta=0{,}01\) se vuelve un arco de \(0{,}02\) metros, y \(J^TJ=\mathrm{diag}(1,4)\) es
la métrica de esa grilla, la \(G\) del cambio anterior. Es la jugada que el profesor llama
``ir a lo chico para volver a lo grande'' (resumen 2025, clase 1).

\subsubsection*{Un mismo objeto con las dos miradas: el cuarto de vuelta}

En el CBC, girar \(90^\circ\) era multiplicar por una matriz:
\[
  R=\begin{pmatrix}0&-1\\1&0\end{pmatrix},\qquad
  R\begin{pmatrix}2\\1\end{pmatrix}=\begin{pmatrix}-1\\2\end{pmatrix}.
\]
Con eso termina la cuenta del CBC. Para el profesor, una rotación implica ``dejar quieto un
punto, mover otro, y en base a esos acomodar todos los demás'' (resumen 2025, clase 1), y
cada propiedad de \(R\) tiene su dibujo.

\emph{Es una operación de un grupo.} Componerla consigo misma da media vuelta,
\(R^2=-I\), y cuatro veces vuelve a la identidad, \(R^4=I\). Su inversa es girar en sentido
contrario, \(R^{-1}=R^3=R^T=\bigl(\begin{smallmatrix}0&1\\-1&0\end{smallmatrix}\bigr)\), y
basta transponer porque \(R\) conserva distancias. En el plano dos giros alrededor del mismo
centro conmutan, porque los ángulos se suman en cualquier orden. Con una afinidad axial ya
no: con \(D=\mathrm{diag}(2,1)\),
\[
  RD=\begin{pmatrix}0&-1\\2&0\end{pmatrix}\neq
  DR=\begin{pmatrix}0&-2\\1&0\end{pmatrix}.
\]

\emph{Se identifica con la trayectoria que genera.} Aplicada repetidas veces, \(R\) lleva
\((2,1)\to(-1,2)\to(-2,-1)\to(1,-2)\to(2,1)\): cuatro puntos sobre la circunferencia de
radio \(\sqrt5\). Si en vez de saltar de a cuarto de vuelta giramos un ángulo \(t\) con
\(R_t=\bigl(\begin{smallmatrix}\cos t&-\sin t\\ \sin t&\cos t\end{smallmatrix}\bigr)\) y
derivamos en \(t=0\), aparece de nuevo \(R\): el profesor anotó que ``la derivada de una
rotación era \(i\)'' (resumen 2025, clase 1). En \((2,1)\) la velocidad es \(R(2,1)^T=(-1,2)\),
perpendicular al radio porque \(2\cdot(-1)+1\cdot2=0\): tangente a la circunferencia. El
cuarto de vuelta es a la vez un salto y el campo de velocidades del giro continuo.

\emph{Tiene un único punto fijo.} \(Rx=x\) pide \((-y,x)=(x,y)\), es decir \(x=-y\) e
\(y=x\), y eso sólo lo cumple el origen. El centro es lo único que no se mueve; si giramos
alrededor de otro centro \(c\), con \(x\mapsto c+R(x-c)\), el punto fijo pasa a ser \(c\).

\emph{No tiene direcciones propias reales.} El polinomio característico es
\[
  \det(R-\lambda I)=\lambda^2+1=0\quad\Longrightarrow\quad\lambda=\pm i .
\]
El hecho geométrico se ve sin el polinomio: para \(v=(x,y)\neq0\), \(v\cdot Rv=-xy+yx=0\),
de modo que cada dirección termina perpendicular a sí misma, y la cuña
\(v\wedge Rv=\det(v\ \ Rv)=x^2+y^2\) no se anula nunca: el detector de colinealidad no
encuentra ninguna dirección que caiga sobre sí misma. Los autovalores complejos son la
manera algebraica de decir que no la hay.

\emph{Es el número \(i\).} Con \(z=2+i\), \(iz=2i+i^2=-1+2i\), que es el punto \((-1,2)\) de
antes; \(i^2=-1\) es \(R^2=-I\), y \(1/i=-i\) es \(R^T\). ``Multiplicar por i es rotar
\(\pi/2\)'' (resumen 2025, clase 1), y la matriz \(R\) es el complejo \(i\) con otro ropaje.

\emph{La escalera.} En \(\R\) los únicos movimientos que conservan distancias y dejan quieto
el \(0\) son \(x\mapsto x\) y \(x\mapsto-x\): un cuarto de vuelta pediría un número \(a\)
con \(a^2=-1\), y no lo hay. En \(\R^2\) aparece, y es \(R\). En \(\R^3\) el cuarto de vuelta
alrededor del eje \(z\) agrega una dirección quieta, el eje: autovalor \(1\) junto a
\(\pm i\). En \(\R^n\) el mismo dibujo se repite: planos que giran como éste y direcciones
que quedan quietas.

\begin{idea}[Qué conviene traer del CBC y qué conviene soltar]
Traemos intactas las cuentas: Gauss, determinantes, inversas, polinomio característico.
Soltamos tres reflejos: leer un punto como su columna de números, leer una matriz como una
tabla y no como un movimiento, y creer que la distancia viene con el espacio. Cada cuenta
vieja responde ahora a una pregunta del terreno: qué se mueve, qué queda quieto, con qué
regla se mide.
\end{idea}

{\small\vspace{1mm}\linkUnidad{00_desaprender_el_cbc}{Simulación interactiva: Unidad 0 — Fundamentos: cuerpos, grupos, puntos, formas lineales y métrica (el idioma de la materia)}\par}

\section{Construcciones, cuentas y ejemplos}

\subsection{Fundamentos: los cinco ladrillos y el idioma de la materia}\label{sec:desaprender}

{\small\vspace{1mm}\linkUnidad{00_desaprender_el_cbc}{Simulación interactiva: Unidad 0 — Fundamentos: cuerpos, grupos, puntos, formas lineales y métrica (el idioma de la materia)}\par}


Desarrollo complementario, armado sobre las clases 1, 4 y 9 y la
bibliografía del campus; no es contenido dictado. El álgebra y el análisis del CBC dejan
reflejos que funcionan en el papel y que esta materia pide cambiar, no porque estén mal
sino porque acá los objetos tienen un significado en el terreno. Para leer el resto con una
sola cabeza alcanzan cinco ladrillos, en orden deductivo: cada uno usa sólo los anteriores.
La Unidad 0 del sitio (\texttt{00\_desaprender\_el\_cbc}, apéndice~\ref{ap:notebooks}) los
trabaja con figuras y ejercicios en Python; la subsección~\ref{sec:puntos-formas-grillas}
los resume en una tabla. Las frases entre comillas son del profesor y están en los
resúmenes 2025 o en las transcripciones 2026; lo demás es desarrollo del apunte.
\begin{enumerate}
  \item \textbf{Conjuntos, operaciones y grupos.} Un grupo es un conjunto con una
  operación interna asociativa, un neutro y una inversa para cada elemento; si el orden no
  importa, conmuta. Es el molde de todo lo que la materia llama \emph{transformación}:
  ``Transformaciones y Operaciones [\ldots] son dos conceptos que van atravesando la
  materia'' (resumen 2025, clase 1), y la inversa, dicha para todos los casos a la vez,
  ``es un elemento de un conjunto que, al operar con otro resulta en el elemento neutro
  respecto de esa operación'' (resumen 2025, clase 10). Los giros de la estación total
  alrededor de su eje vertical conmutan; los de \(\R^3\) no, ``porque la composición de
  transformaciones lineales no es conmutativa'' (resumen 2025, clase 1). En NumPy:
  \texttt{A @ B}, \texttt{np.eye(n)}, \texttt{np.linalg.solve(A, y)}.
  \item \textbf{Los cuerpos \(\R\) y \(\mathbb{C}\).} Un cuerpo tiene dos operaciones, suma
  y producto, cada una un grupo conmutativo (el producto, sin el \(0\)), y el producto
  distribuye. Un real escala la recta; agregar \(i^2=-1\) equivale exactamente a pasar a las
  matrices \(\bigl(\begin{smallmatrix}a&-b\\ b&a\end{smallmatrix}\bigr)\), que respetan
  suma, producto e inverso y rotan \(\arg z\) mientras escalan \(|z|\): ``la matriz de
  rotación de 2x2 que vimos en algebra del CBC hace exactamente la misma operación que
  hacíamos con los complejos'' (resumen 2025, clase 1). La ``interesantísima equivalencia
  en la estructura de los numeros complejos y algunas matrices de 2x2'' el profesor la
  remite a Gentile, \emph{Complejos y Generalización} (bibliografía del campus,
  apéndice~\ref{ap:bibliografia}). En NumPy:
  \texttt{np.array([[z.real, -z.imag], [z.imag, z.real]])}.
  \item \textbf{Espacio afín y espacio vectorial: todo vector es un punto.} Un espacio
  vectorial es un conjunto con suma y producto por escalar que respetan linealidad (``un
  conjunto al que se asocian dos operaciones [\ldots] que cumplen con algunas propiedades
  básicas a las que llamamos `linealidad'\,'', resumen 2025, clase 4); de ahí salen
  combinación lineal, independencia, base y dimensión. Un espacio afín es un conjunto de
  puntos sin origen: la diferencia de dos puntos es un vector. El terreno es afín (hay
  mojones \(P\), \(Q\), no flechas); \(V\) nace cuando plantamos la estación total en un mojón
  \(O\) y cada mojón pasa a ser \(P-O\), o cuando homogeneizamos \((x,y)\mapsto(x,y,1)\). Toda
  transformación lineal ``conserva el cero'' (19/8/2026), por eso las traslaciones piden
  homogéneas. En el espíritu de lo que dice el profesor sobre los campos vectoriales
  (``tiene sentido pensarlos como puntos'', 14/10/2026, transcripción auto-traducida): en
  esta materia un vector es un mojón visto desde \(O\) o un desplazamiento entre dos. En
  NumPy, punto y desplazamiento son el mismo \texttt{np.array([x, y])}; independencia:
  \texttt{np.linalg.matrix\_rank}.
  \item \textbf{El dual \(V^{*}\): una fila es una forma lineal, y el gradiente también.}
  Una forma lineal es una función lineal \(\omega:V\to\R\); las formas de \(V\) forman un
  espacio vectorial de la misma dimensión, con base las selectoras de coordenadas (``las
  más básicas de esas funciones son las funciones selectoras de coordenadas'', resumen
  2025, clase 4). Se escribe como fila y evalúa fila por columna, sin producto interno; se
  ve como una regla graduada, un mazo de rectas de nivel equiespaciadas. ``No existe el
  vector fila. Lo que comúnmente llamamos vector fila se llama forma lineal''
  (30/9/2026, transcripción auto-traducida). Si los puntos van con \(A\), la regla se
  arrastra con \(A^{-1}\) para que la lectura no cambie: ``una forma lineal se transforma
  con una transformación lineal del lado `contrario' que un vector'' (resumen 2025,
  clase 4). Y \(df=\nabla f\,dx\) es una de ellas: ``dato clave: el gradiente es una forma
  lineal no un vector'' (resumen 2025, clase 9). En NumPy: \texttt{w @ x} lee,
  \texttt{np.linalg.solve(A.T, w)} arrastra.
  \item \textbf{Producto interno y métrica \(G\): la cinta métrica.} Un producto interno
  es \(\langle u,v\rangle_G=u^TGv\) con \(G\) simétrica definida positiva; recién con él
  hay largos y ángulos, y es el único puente entre \(V\) y \(V^{*}\): baja un vector a forma
  (\(w=Gv\)) y sube una forma a vector (\(v=G^{-1}w\)). ``Lo único que define a la base que
  llamamos canónica es que sobre ella tenemos definida la noción de distancia a través del
  producto interno'' (resumen 2025, clase 6): en esa base \(G=I\) y la cinta coincide con la
  grilla; en la grilla de un plano estirado, \(G'=A^{-T}A^{-1}\); en \((\varphi,\lambda)\)
  sobre una superficie, \(G\) es la primera forma fundamental, que cambia de punto a punto.
  En NumPy: \texttt{a @ g @ b}, \texttt{np.linalg.solve(g, w)}, \texttt{J.T @ J}.
\end{enumerate}

El puente al cálculo en superficies, que es lo que hace de las clases 9 a 14 la misma
materia que las clases 1 a 8: derivar una parametrización es usar una matriz lineal, la
jacobiana \(J\), en el espacio tangente (ladrillos 1 y 3); el gradiente es una forma del
dual (ladrillo 4); la primera forma fundamental \(J^TJ\) es la métrica del tangente
(ladrillo 5); y curvarse, hessiana o segunda forma, es una forma cuadrática: un número por
cada dirección. La subsección~\ref{sec:puntos-formas-grillas} pone los tres objetos en
una tabla y la Unidad 0 los manipula.

\begin{resultado}[Las tres ecuaciones maestras de la materia (\(Ax=b\), \(Ax=\lambda x\), \(Ax\approx b\))]
Fijados los cinco ladrillos, casi todas las cuentas lineales del curso responden a una de tres preguntas (las que quedan afuera son las de medir con la métrica, \(J\) y \(J^TJ\), en las unidades 02, 10, 13 y 14):

\smallskip
\small
\begin{tabularx}{\linewidth}{>{\raggedright\arraybackslash}p{19mm} >{\raggedright\arraybackslash}X >{\raggedright\arraybackslash}p{47mm}}
\textbf{Ecuación} & \textbf{Pregunta de campo} & \textbf{Dónde la desarrollamos} \\
\hline
\(Ax=b\) &
\textbf{Transformar o deshacer.} Evaluar \(Ax\) es transformar: conocemos \(A\) y calculamos a dónde va un mojón. Resolver \(Ax=b\) es la pregunta inversa: conocemos \(A\) y el resultado \(b\) y buscamos de dónde vino, \(x=A^{-1}b\); leer un punto en otra base es eso mismo (las coordenadas nuevas son la \(x\) que cumple \(Px=b\)). Hallar la homografía que rectifica los cuatro puntos de una foto de dron también es un \(Ax=b\), pero con la matriz como incógnita: un sistema lineal \(8\times8\) en los coeficientes de \(H\). &
\linkUnidad{01a_punto_quieto}{01a}, \linkUnidad{01b_rotacion_plano}{01b}, \linkUnidad{01c_rotaciones_3d}{01c}, \linkUnidad{03_reflexiones_similaridades}{03}, \linkUnidad{04_afines_homogeneas}{04}, \linkUnidad{05_proyectivas}{05}, \linkUnidad{06_formas_lineales}{06}, \linkUnidad{08_cambio_de_base}{08} \\
\hline
\(Ax=\lambda x\) y \(Av_i=\sigma_i u_i\) &
\textbf{Buscar direcciones privilegiadas.} Las que no giran (autovectores): el eje de giro en \(\mathbb{R}^3\), los puntos fijos de una afinidad o proyectiva, las direcciones principales de curvatura del elipsoide, cuyas curvaturas principales \(1/M\) y \(1/N\) son los autovalores de \(I^{-1}II\) (el operador de forma). Y, con la SVD, el par de bases perpendiculares que sigue perpendicular: los \(v_i\) (direcciones propias de \(A^TA\)) van a los \(u_i\) girando, pero sin perder el ángulo recto; son los semiejes de la elipse de Tissot. &
preparación: \linkUnidad{07a_cuna_determinante}{07a}; \linkUnidad{07b_nucleo_autovectores}{07b}, \linkUnidad{07c_autovalores_3d_puntos_fijos}{07c}, \linkUnidad{08_cambio_de_base}{08}, \linkUnidad{09_cambio_de_base_svd}{09}, \linkUnidad{12_curvatura}{12}, \linkUnidad{13_elipsoide}{13}, \linkUnidad{15_transformaciones_conformes}{15} \\
\hline
\(Ax\approx b\) &
\textbf{Aproximar cuando sobran mediciones y hay error de campaña.} Hay más ecuaciones que incógnitas y el vector medido \(b\) cae fuera de la imagen de \(A\); buscamos \(\hat{x}=A^{+}b\), cuya imagen \(A\hat{x}\) es la proyección perpendicular de \(b\) sobre el subespacio columna \(\mathrm{Im}(A)\): regresión e intersección de círculos por distancias. Procrustes y regresión total también son mínimos cuadrados, con otra geometría (restricción ortogonal; distancia perpendicular), y se resuelven con la SVD de una matriz armada con los datos. &
\linkUnidad{02_producto_interno}{02}, \linkUnidad{09_cambio_de_base_svd}{09}, \linkUnidad{11_cuatro_problemas}{11} \\
\end{tabularx}
\normalsize

\smallskip
Los cinco ladrillos son el idioma de la materia, y estas tres ecuaciones son las tres frases que se dicen con él: transformamos o deshacemos, buscamos direcciones privilegiadas, o proyectamos lo medido sobre lo que el modelo puede explicar. Cada vez que aparece una cuenta nueva, nos preguntamos cuál de las tres es.
\end{resultado}

\subsection{Convención para componer}

En este apunte, el producto \(F G\) significa ``primero aplicar \(G\) y después \(F\)'':
\[
  (F G)(x)=F\bigl(G(x)\bigr).
\]
Esta convención es importante porque, en general, cambiar el orden cambia el resultado.

\section{En profundidad}

\subsection{La mirada geométrica}

El hilo conductor de la materia no es memorizar matrices, sino aprender a \emph{ver} qué
transformación describe cada fórmula. Antes de calcular conviene interrogar la figura.

\begin{center}
\begin{tabularx}{\textwidth}{>{\bfseries}p{38mm} X}
\toprule
Pregunta geométrica & Qué revela\\
\midrule
¿Qué queda fijo? & centro, eje, recta o dirección privilegiada\\
¿Cómo se mueve un punto? & trayectoria recta, circular o curva\\
¿Qué se conserva? & distancia, ángulo, paralelismo, razón, área u orientación\\
¿Qué cambia? & escala, dirección, forma, curvatura o sistema de referencia\\
¿Qué pasa con una figura simple? & la imagen de un segmento, triángulo o paralelogramo
determina la transformación\\
¿Qué hace la composición? & permite construir transformaciones complejas a partir de
operaciones simples\\
\bottomrule
\end{tabularx}
\end{center}

\begin{idea}[Regla de lectura]
Para cada teorema se seguirá este orden:
\begin{enumerate}
  \item \textbf{Construcción}: dibujar el punto y su imagen.
  \item \textbf{Movimiento}: describir la trayectoria y el sentido.
  \item \textbf{Invariantes}: identificar qué no cambia.
  \item \textbf{Composición}: seguir los pasos intermedios en el orden correcto.
  \item \textbf{Traducción algebraica}: escribir la matriz o fórmula que resume el dibujo.
  \item \textbf{Verificación}: usar un ejemplo numérico para comprobar, no para sustituir,
  la comprensión geométrica.
\end{enumerate}
\end{idea}

\begin{resultado}[La matriz como geometría comprimida]
Las columnas de una matriz \(A\) son las imágenes de los vectores de la base:
\[
  A=
  \begin{pmatrix}\vert&\vert\\Ae_1&Ae_2\\\vert&\vert\end{pmatrix}.
\]
Por eso, para interpretar geométricamente una transformación lineal basta observar qué hace
con dos direcciones independientes y con el paralelogramo que generan.
\end{resultado}


\subsubsection{Los tres objetos que no hay que confundir: puntos, reglas de medir y grillas}\label{sec:puntos-formas-grillas}

{\small\vspace{1mm}\linkAnimacion{01_dualidad_gradiente}{Animación: puntos con \(A\), reglas de medir con \(A^{-1}\)}\par}

Síntesis del curso, como brújula para el resto del apunte: en la materia circulan tres objetos
que se escriben con números parecidos y son cosas distintas. El hilo conductor es que
todo vector es un mojón o un desplazamiento entre mojones, y toda matriz es una manera de
mover el espacio o de cambiar la grilla con que se lo mide.

Los dos ladrillos de abajo de esa tabla, desarrollados en la subsección~\ref{sec:desaprender}
y en la Unidad 0, son el idioma en que están escritas las tres filas:

\begin{description}
  \item[Grupos de transformaciones (ladrillo 1).] Las transformaciones de un mismo tipo
  (rotaciones, traslaciones, afinidades, homografías) forman un grupo: componer dos da otra
  del mismo tipo, la identidad es el neutro y cada una tiene inversa. ``Transformaciones y
  Operaciones [\ldots] son dos conceptos que van atravesando la materia'' (resumen 2025,
  clase 1); la inversa ``es un elemento de un conjunto que, al operar con otro resulta en el
  elemento neutro respecto de esa operación'' (resumen 2025, clase 10). Cuándo conmuta el
  grupo es la pregunta que decide si se puede armar un sistema de coordenadas: las
  rotaciones del plano y las traslaciones conmutan; las rotaciones de \(\R^3\) no
  (sección~\ref{sec:orden-importa}). En NumPy la operación es \texttt{A @ B}, el neutro
  \texttt{np.eye(n)} y la inversa aplicada a un dato \texttt{np.linalg.solve(A, y)}.
  \item[Los cuerpos \(\R\) y \(\mathbb{C}\) como matrices (ladrillo 2).] Los escalares
  escalan: un real \(a\) actúa sobre la recta como \(x\mapsto ax\), con inversa \(1/a\) y
  determinante \(|a|\). Un complejo \(z=a+ib\) actúa sobre el plano como la matriz
  \[
    M(z)=\begin{pmatrix}a&-b\\ b&a\end{pmatrix},\qquad
    M(z_1+z_2)=M(z_1)+M(z_2),\quad M(z_1z_2)=M(z_1)M(z_2),\quad M(1/z)=M(z)^{-1},
  \]
  es decir rota \(\arg z\) y escala \(|z|\) a la vez, con \(\det M(z)=|z|^2\): las matrices
  de esa forma \emph{son} \(\mathbb{C}\) con otro ropaje. ``La matriz de rotación de 2x2 que
  vimos en algebra del CBC hace exactamente la misma operación que hacíamos con los
  complejos'' (resumen 2025, clase 1; la equivalencia de estructuras la remite a Gentile,
  \emph{Complejos y Generalización}, apéndice~\ref{ap:bibliografia}). Sección~\ref{sec:rot-complejos}.
\end{description}

Con el grupo y el cuerpo puestos, los tres objetos:

{\small
\begin{tabularx}{\textwidth}{>{\bfseries\raggedright\arraybackslash}p{27mm}>{\raggedright\arraybackslash}p{30mm}YY}
\toprule
Objeto & Cómo se escribe, dónde vive & Qué es geométricamente & Cómo se transforma\\
\midrule
Punto / vector & columna \(v\); vive en el espacio \(V\) &
una flecha: un mojón (desde el origen) o un desplazamiento entre dos mojones &
con la matriz, por la izquierda: \(v\mapsto Av\)\\
Forma lineal / gradiente & fila \(\omega\); vive en el dual \(V^{*}\) &
un mazo de rectas (en \(\R^3\), planos) de nivel equiespaciadas: una regla graduada
tirada en el plano; evalúa un vector \emph{sin} producto interno, mirando qué recta del
mazo pasa por él &
del lado contrario: \(\omega\mapsto\omega A^{-1}\), para que \(\omega v\) no cambie\\
Producto interno / métrica & matriz simétrica \(G\) (en una superficie, la primera forma
fundamental \(E,F,G\)) &
la grilla con que se mide: permite largos y ángulos, \(\|v\|^2=v^TGv\); es el único
puente entre una forma lineal y una dirección perpendicular (``el gradiente es
perpendicular a la curva de nivel'' necesita a \(G\)) &
como forma cuadrática: \(G\mapsto J^TGJ\) (sección~\ref{sec:metrica})\\
\bottomrule
\end{tabularx}
}

Sin producto interno no hay ``vector perpendicular a una recta'': hay la recta (forma) y
hay los puntos (vectores), y nada más. Con producto interno, la forma \(\omega\) se
``disfraza'' del vector \(G^{-1}\omega^T\) y recién ahí se puede hablar de la dirección
normal. En la base canónica \(G=I\) y el disfraz es transponer la fila, por eso en el
plano la confusión no cuesta nada; en una superficie, donde \(E,F,G\) cambian de punto a
punto, cuesta metros.

\begin{atencion}[Por qué el gradiente no es un vector]
\(df=\nabla f\,dx\) es una forma lineal: toma un desplazamiento \(dx\) y devuelve cuánto
cambia \(f\); sus curvas de nivel son el plano tangente a la gráfica, y su recta de nivel
cero es la tangente a la curva de nivel de \(f\). Prueba de que no es un vector: si
estiramos el eje \(x\) al doble, los puntos duplican su primera coordenada, pero las marcas
de la regla \(\nabla f\) a lo largo de \(x\) se separan a la mitad (hay que avanzar el
doble para que \(f\) suba lo mismo), es decir la componente se divide por dos: se
transforma al revés que un vector. El profesor se corrigió en vivo el 14 de octubre de
2026, mientras resolvía el plano normal a una curva intersección de dos superficies: ``el
gradiente de una función es un vector\ldots{} no, mentí. Lo dije mal. [\ldots] El gradiente
de la función es una forma lineal, no un vector, porque me permite hacer una aproximación
lineal de la función'' (transcripción auto-traducida, \file{fuentes/transcripciones/2026-10-14.md}).
Y es lo que el resumen 2025 de la clase 9 llama ``dato clave: el gradiente es una forma
lineal no un vector''. Lo que sí es un vector es la normal, que se obtiene con el producto
vectorial de dos tangentes, o levantando el gradiente con la métrica.
\end{atencion}

\begin{resultado}[El puente al cálculo en superficies]
Los cinco ladrillos de la subsección~\ref{sec:desaprender} son también el idioma de las
clases 9 a 14. Derivar una parametrización \(r(u,v)\) es usar una matriz lineal, la
jacobiana \(J=(r_u\ \ r_v)\), en el espacio tangente: entra un pasito \((du,dv)\) de la
grilla, sale el pasito \(J\,(du,dv)^T\) sobre la superficie (ladrillos 1 y 3,
sección~\ref{sec:variedades}). El gradiente de una función escalar es una forma del dual,
\(df=\nabla f\,dx\) (ladrillo 4). La primera forma fundamental \(J^TJ=
\bigl(\begin{smallmatrix}E&F\\ F&G\end{smallmatrix}\bigr)\) es la métrica del tangente,
la única cinta con que se miden largos y ángulos en la grilla \((u,v)\) (ladrillo 5,
sección~\ref{sec:metrica}); por eso la perpendicular dibujada en \((\varphi,\lambda)\) no
es la del terreno, y el gradiente levantado con esa métrica, \(G^{-1}\nabla f^{T}\), sí lo
es. Y curvarse, la hessiana o la segunda forma fundamental, es una forma cuadrática: un
número por cada dirección, cuya diagonalización da las direcciones principales
(sección~\ref{sec:curvatura}).
\end{resultado}
